\NoteChapter{4}{Chapter 5 Hilbert空间}
\NoteSection{1}{Hilbert空间的基本性质}


\begin{proposition}
    Hilbert空间是一致凸空间，
    因而是自反空间.
\end{proposition}

\begin{proof}
    对任意的$\varepsilon >0$，
    $u,v$满足$\| u \| \le 1, \|v\| \le 1, \| u -v\| > \varepsilon$.
    由平行四边形恒等式

    \[
    \left\| \frac{u+v}{2} \right\|^2 = \frac{1}{2}(\|u\|^2 + \| v \|^2) - 
    \left\| \frac{u-v}{2} \right\|^2 < 1 - \frac{\varepsilon^2}{4}
    \]
\end{proof}

\begin{theorem}[点到闭凸集上的投影]
    设$K \subset H$是非空闭凸集，
    则对任意的$f \in H$，
    存在唯一的$u \in K$，
    使得$\|f-u\|=d(f , K)$.
    记$u = P_K f$.
    这样的$u$是$K$中唯一满足如下条件的元素：
    $$
    (f-u,v-u) \le 0, \quad v \in K
    $$
\end{theorem}

\begin{remarkbox}
    \textbf{Remark}：
    对于存在性，
    有一个简短的证明：
    考虑函数$\varphi(v) = \| f-v \|$.
    它是连续的凸函数，
    当$\left\| v \right\| \to \infty$时，
    $\varphi(v) \to \infty$.
    且$H$是自反空间.
    因此$\varphi$能在$K$上取到最小值.
    接下来给出一个更直接的证明.
\end{remarkbox}

\begin{proof}~

    (a) 存在性：取一列$v_n \in K$，
    使得$d_n = \left\| f-v_n \right\| \to d = \inf_{v \in K} \left\| f - v\right\| $.
    只需证明$v_n$是柯西列.
    利用平行四边形恒等式去估计$\left\| v_n - v_m \right\| = \left\| (f - v_n) - (f - v_m) \right\|$即可.

    \smallskip

    (b) 唯一性：
    只需证明不等式$(f-u , v-u) \le 0$，
    唯一性便是显然的.

    取$w \in H$，
    使得对充分小的正数$\varepsilon$，
    满足$u + \varepsilon w = v\in K$.
    则$\left\| f - u\right\| \le \left\| f - u - \varepsilon w\right\|$.
    两边取平方，有
    \[
    \left\| f-u  \right\|^2 \le \left\| f-u\right\|^2 - 2 (f-u , \varepsilon w) + \left\| \varepsilon w \right\|^2 
    \]

    即
    \[
    (f-u , w) \le \varepsilon \left\| w \right\|^2
    \]

    取$\varepsilon \to 0$，则
    \[
    (f-u , w) \le 0
    \]

    这等价于对任意的$v \in K$，
    \[
    (f-u , \varepsilon w) = (f-u , v - u) \le 0
    \]
\end{proof}

\begin{corollary}
    设$M \subset H$是闭子空间，
    $f \in H$，
    则对于$u \in M$，
    $u = P_M f$当且仅当$(f-u , v) = 0$对任意$v \in M$成立.
    此时$P_M$是线性算子，
    称为正交投影算子.
\end{corollary}

\NoteSection{2}{Gelfand三元组}


\begin{theorem}[Riesz表示定理]
    对任意的$\varphi \in H^\star$，
    存在唯一的$f \in H$，
    使得$\varphi (u) = (f,u) , \forall \in H$.
    且$\| f \| = \| \varphi \|_{H^\star}$.
\end{theorem}

Riesz定理表明，
$H$与$H^{\star}$之间存在一个自然的等距同构，
从而我们可以把$H$与$H^{\star}$看成完全等同的空间
（将$H \to H^\star$的等距同构视为恒等映射）.
但要注意的是，
这种等同只是在内积结构的意义上等同，
而不是在集合元素的意义上等同.
接下来我们通过Gelfand三元组的构造说明这一点，
并给出一个具体例子.

\smallskip

设$H$是Hilbert空间，
Banach空间$V$可连续地嵌入$H$，
从而可以理解为$V \subset H$.

自然地，我们可以将$H^\star$中的元素（线性泛函）限制在$V$上，
成为$V^\star$中的元素.
记这个限制映射为$T: H^\star \to V^\star$.
具体来说，$T$满足

\[
\langle T\varphi , v \rangle_{V} = 
\langle \varphi , v \rangle_{H} 
\]


一般来说，$T$不是满射.
$T$给出了$H^\star \to V^\star$的一个嵌入.
如果把$H$与$H^\star$视为等同，
我们有

\[
V \subset H \cong H^\star \subset V^{\star}
\]

这就是Gelfand三元组.
此时，如果我们再假设$V$是Hilbert空间，
其上定义的内积$(\cdot , \cdot)_V$与$H$的内积不一定相同.
我们发现不能再将$V$与$V^\star$按$V$上的内积视为等同了，
这将与$H = H^\star$矛盾.
当然，$V \to V^\star$的等距同构依然是存在的，
只不过不能把它看成$V$上的恒等映射了.

\smallskip

一个具体的例子如下.
设$H$是平方可和序列空间$\ell^2$，
装备了$\ell^2$的标准内积.
设$V \subset H$是加权平方可和空间：

\[
V = \left\{ u = (u_n) , \sum_{n=1}^{\infty} n^2 u_n^2 < \infty \right\}
\]

其上的内积定义为

\[
(u , v)_{V} = \sum_{n=1}^\infty n^2 u_n v_n
\]


根据Riesz表示定理，
对任意的$F \in V^{\star}$，
存在唯一的$u \in V$，
使得$\langle F , v \rangle_V = (u , v)_V = \sum_{n=1}^\infty n^2 u_n v_n$.
特别地，若
$F = T \varphi$，
其中$ \varphi \in H^\star = H$，
考虑到$T$所满足的关系

\[
\langle T\varphi , v \rangle_{V} = 
\langle \varphi , v \rangle_{H} = 
 \sum_{n=1}^\infty n^2 u_n v_n
\]

$\varphi$作为$H$中元素，
应为$n^2 u_n$.
则有$\dfrac{\varphi_n^2}{n^2} = n^2 u_n^2$是可和序列.
因此，$T: u_n \mapsto n^2 u_n$是$V \to V^\star$的Riesz同构映射.
$V^\star$应自然地认作

\[
V^\star = 
\left\{ f = (f_n) , \quad \sum_{n=1}^{\infty} \frac{1}{n^2} f_n^2 < \infty \right\}
\]


显然，$V$是$V^\star$的真子集.

\NoteSection{3}{Lax-Milgram定理}

\begin{definition}
    设有双线性形式$a : H \times H \to \mathbb{R}$，

    (i) 若存在常数$C$，
    使得$|a (u,v) | \le C |u||v|, \quad u , v \in H$
    则称$a$是连续的.

    (ii) 若存在常数$\alpha$，
    使得$a(v,v) \ge \alpha |v|^2, \quad v \in H$，
    则称$a$是强迫的. 
\end{definition}

\begin{remarkbox}
    \textbf{Remark}：
    连续性控制了双线性形式的上界，
    强迫性控制了双线性形式的下界.
    既连续又强迫的双线性形式可以看成内积的推广.
    内积本身也是一种既连续又强迫的双线性形式.
\end{remarkbox}

\begin{theorem}[Stampacchia定理]
    设$a$是$H$上的连续强迫双线性形式，
    $K \subset H$是非空闭凸集.
    则对任意的$\varphi \in H^\star$，
    存在唯一的$u \in K$，
    使得$a(u , v-u) \ge \varphi (v - u)$.

    若$a$还是对称的，
    则$u$满足上述性质当且仅当$u \in K$且：

    \[
    \frac{1}{2} a(u,u) - \varphi(u) = \min_{v \in K} \left\{ \frac{1}{2} a(v,v) - \varphi (v) \right\}
    \]
\end{theorem}

